%!TEX root = 1_a_Generalised IIT.tex
\section{Summary \& Outlook} \label{sec:discussion}

In this article, we have propounded the mathematical structure of Integrated Information Theory.
First, we have studied which exact structures the IIT algorithm uses in the mathematical description of physical systems, on the one hand, and in the mathematical description of conscious experience, on the other. Our findings are the basis of definitions of a physical system class $\Sys$ and a class $\Expcat$ of experience spaces, and allowed us to view IIT as a map $\Sys \rightarrow \Expcat$.

Next, we needed to disentangle the essential mathematics of the theory from auxiliary formal tools used in the contemporary definition. To this end, we have introduced the precise notion of decomposition of elements of an experience space required by the IIT algorithm. The pivotal cause-effect repertoires are examples of decompositions so defined, which allowed us to view any particular choice, e.g. the one of `classical' IIT developed by Tononi et. al., or the one of `quantum' IIT recently introduced by Zanardi et. al. as data provided to a general IIT algorithm.

The formalization of cause-effect repertoires in terms of decompositions then led us to define the essential ingredients of IIT's algorithm concisely in terms of integration levels, integration scalings and cores. These definitions describe and unify recurrent mathematical operations in the contemporary presentation, and finally allowed to define IIT completely in terms of a few lines of definition.

Throughout the paper, we have taken great care to make sure our definitions reproduce exactly the contemporary version of IIT 3.0. The result of our work is a mathematically rigorous and general definition of Integrated Information Theory. This definition can be applied to any meaningful notion of systems and cause-effect repertoires, and we have shown that this allows to overcome most of the mathematical problems of the contemporary definition identified to date in the literature.

We believe that our mathematical reconstruction of the theory can be the basis for refined mathematical and philosophical analysis of IIT. We also hope that this mathematisation may make the theory more amenable to study by mathematicians, physicists, computer scientists and other researchers with a strongly formal background.


\subsection{Process Theories}
Our generalization of IIT is axiomatic in the sense that we have only included those formal structures in the definition which are necessary for the IIT algorithm to be applied. This ensured that our reconstruction is as general as possible, while still true to IIT 3.0. As a result, several notions used in classical IIT, e.g., system decomposition, subsystems or causation, are merely defined abstractly at first, without any reference to the usual interpretation of these concepts in physics.

In the related article~\cite{IITProcess}, we show that these concepts can be meaningfully defined in any suitable \emph{process theory} of physics, formulated in the language of \emph{symmetric monoidal categories}. This approach can describe both classical and quantum IIT and yields a complete formulation of contemporary IIT in a categorical framework. 

\subsection{Further Development of IIT}
IIT is constantly under development, with new and refined definitions being added every few years. We hope that our mathematical analysis of the theory might help to contribute to this development. E.g., the working hypothesis that IIT is a fundamental theory, i.e. describes reality as it is, implies that technical problems of the theory need to be resolved. We have shown that our formalization allows address the technical problems mentioned in the literature. However, there are others which we have not addressed in this paper.

Most crucially, the IIT algorithm uses a series of maximalization and minimalization operations, unified in the notion of \emph{core} subsystems in our formalization. In general, there is no guarantee that these operations lead to unique results, neither in classical nor quantum IIT. Using different cores has major impact on the output of the algorithm, including the $\Phi$ value, which is a case of ill-definedness.\footnote{The problem of `unique existence' has been studied extensively in category theory using \emph{universal properties} and the notion of a \emph{limit}. Rather than requiring that each $E \in \Exp$ come with a metric, it may be possible to alter the IIT algorithm into a well-defined categorical form involving limits to resolve this problem.}

Furthermore, the contemporary definition of IIT as well as our formalization rely on there being a finite number of subsystems of each system, which might not be the case in reality. Our formalisation may be extendable to the infinite case by assuming that every system has a fixed but potentially infinite indexing set 
$\Sub(S)$, so that each $\Sub_s(S)$ is the image of a mapping $\Sub(S) \times \St(S) \to \Sys$, but we have not considered this in detail in this paper.

Finally, concerning more operational questions, it would be desirably to develop the connection to empirical measures such as the Perturbational Complexity Index PCI \cite{casarotto2016stratification,casali2013theoretically} in more detail, as well as to define a controlled approximation of the theory whose calculation is less expensive. Both of these tasks may be achievable by substituting parts of our formalization with simpler mathematical structure.
%as well as the stability w.r.t. perturbations of a system's state or time evolution

On the conceptual side of things, it would be desirable to have a more proper understanding of how the mathematical structure of experiences spaces corresponds to the phenomenology of experience, both for the general definition used in our formalization and the specific definitions used in classical and quantum IIT. In particular, it would be desirable to understand how it relates to the important notion of qualia, which is often asserted to have characteristic features such as ineffability, intrinsicality, non-contextuality, transparency or homogeneity~\cite{metzinger2006grundkurs}. For a first analysis towards this goal, cf.~\cite{kleiner2019models}.



